\section*{Chapter \ref{ch:g12:solow-first-look} --- The Solow Model: A First Look}
% figdata: figdata/grade-12/05-solow-first-look/out/ (conv, recovery, paths, solow_key); Brenland invented.

\begin{solution}{exo:g12:solow-first-look:1}
$k_1 = 0.2 \times 1 + 0.9 \times 1 = 1.1$; $k_2 = 0.2\sqrt{1.1} + 0.9 \times 1.1 \approx 0.210 + 0.990 = 1.200$.
\end{solution}

\begin{solution}{exo:g12:solow-first-look:2}
$k_1 = 0.2 \times 3 + 0.9 \times 9 = 8.7$; $k_2 = 0.2\sqrt{8.7} + 0.9 \times 8.7 \approx 0.590 + 7.830 = 8.420$.
\end{solution}

\begin{solution}{exo:g12:solow-first-look:3}
$k^* = (0.2/0.1)^{1/0.5} = 4$; $y^* = \sqrt{4} = 2$.
\end{solution}

\begin{solution}{exo:g12:solow-first-look:4}
$k^* = (0.6/0.06)^{3/2} = 10^{1.5} \approx 31.6$; $y^* = 2 \times 10^{0.5} \approx 6.32$.
\end{solution}

\begin{solution}{exo:g12:solow-first-look:5}
Korea: about 1400 dollars in 1960 and 5.8\,\% a year. Ireland: about 7700 dollars and 3.8\,\% a year. No: some of the poorest countries grew slowly or shrank (the
Democratic Republic of the Congo by $-1.75\,\%$ a year), while some richer ones grew fast.
\end{solution}

\begin{solution}{exo:g12:solow-first-look:6}
$k^*$ is multiplied by $2^{1/(1 - \alpha)} = 2^{1.5} \approx 2.83$, and $y^*$ by $2^{\alpha/(1 - \alpha)} = 2^{0.5} \approx 1.41$.
\end{solution}

\begin{solution}{exo:g12:solow-first-look:7}
$k^*$ by $1.5^{2} = 2.25$, $y^*$ by $1.5^{1} = 1.5$.
\end{solution}

\begin{solution}{exo:g12:solow-first-look:8}
At $k = 1$: $0.2 - 0.1 = 10\,\%$. At $k = 2$: $0.2/\sqrt{2} - 0.1 \approx 4.1\,\%$. At $k = 4$: $0.1 - 0.1 = 0$.
\end{solution}

\begin{solution}{exo:g12:solow-first-look:9}
The gap is 3, so half is closed when $k \ge 2.5$: $k_{17} \approx 2.486$ and $k_{18} \approx 2.553$, so 18 periods.
\end{solution}

\begin{solution}{exo:g12:solow-first-look:10}
About 161 from 100 and 528 from 600, after 10 years; the \omterm{def:g12:solow-first-look:steady-state}{steady state} is 358.4.
\end{solution}

\begin{solution}{exo:g12:solow-first-look:11}
At a \omterm{def:g12:solow-first-look:steady-state}{steady state} $k_{t+1} = k_t = k$: $(1 + n)k = sAk^{\alpha} + (1 - \delta)k$, so $sAk^{\alpha} = (\delta + n)k$. Then $k^* = (sA/(\delta + n))^{1/(1 - \alpha)} =
(0.2/0.12)^2 \approx 2.78$, lower than 4: more workers dilute the capital.
\end{solution}

\begin{solution}{exo:g12:solow-first-look:12}
Across the 111 countries, the correlation between income in 1960 and growth after it is about zero ($-0.08$): poor countries did not, as a group, grow faster. The model
predicts faster growth only for a country below its own \omterm{def:g12:solow-first-look:steady-state}{steady state}; poor countries differ in saving, schooling, technology and institutions, and so in \omterm{def:g12:solow-first-look:steady-state}{steady states}.
Among twenty similar countries, where \omterm{def:g12:solow-first-look:steady-state}{steady states} are closer, the correlation is $-0.66$.
\end{solution}

\begin{solution}{exo:g12:solow-first-look:13}
Since $\alpha - 1 < 0$, $k^{\alpha - 1}$ decreases as $k$ rises, and so does $sAk^{\alpha - 1} - \delta$; it is zero where $sAk^{\alpha} = \delta k$, at $k^*$. By the
proposition, $k_t \to k^*$, so the growth rate of $k$ tends to 0, and output per worker $Ak_t^{\alpha}$ tends to the constant $A(k^*)^{\alpha}$. A higher \omterm{def:g12:solow-first-look:saving-rate}{saving rate} gives a
higher $k^*$, but the growth rate again tends to 0.
\end{solution}

\begin{solution}{exo:g12:solow-first-look:14}
$k$ falls from 4 to 2, $y$ from 2 to $\sqrt{2} \approx 1.41$: by about 29\,\%. Then $k_1 = 0.2\sqrt{2} + 0.9 \times 2 \approx 2.083$, growth of about 4.1\,\%, against 0 at the
\omterm{def:g12:solow-first-look:steady-state}{steady state}. Germany's output per person in 1946 was less than half that of 1938, and it grew much faster than the United States in the years after.
\end{solution}

\begin{solution}{exo:g12:solow-first-look:15}
$k^*$ is proportional to $A^{1/(1 - \alpha)} = A^{1.5}$: a ratio of $2^{1.5} \approx 2.83$. $y^* = A(k^*)^{1/3}$ is proportional to $A^{1.5}$ too: also 2.83. The country with
twice the technology also has 2.83 times the capital, because it saves out of a larger output: part of the capital gap is itself a consequence of the technology gap.
\end{solution}

\begin{solution}{pb:g12:solow-first-look:1}
\textbf{1.} $88\,484/393\,264 \approx 0.225$.

\textbf{2.} $218.42^{0.35} \approx 6.588$, so $A = 66.97/6.588 \approx 10.166$.

\textbf{3.} $k^* = (0.225 \times 10.166/0.05)^{1/0.65} = 45.75^{1.538} \approx 358.4$.

\textbf{4.} $y^* = 10.166 \times 358.4^{0.35} \approx 79.6$ thousand dalens.

\textbf{5.} $218.42/358.4 \approx 61\,\%$ of its steady-state capital, and $66.97/79.6 \approx 84\,\%$ of its steady-state output, per worker.

\textbf{6.} $0.225 \times 10.166 \times 6.588 + 0.95 \times 218.42 \approx 15.07 + 207.50 = 222.57$.

\textbf{7.} $222.57/218.42 - 1 \approx 1.9\,\%$.

\textbf{8.} Saving per worker relative to capital, $sAk^{\alpha - 1}$, is highest when $k$ is low, while depreciation is a fixed share of $k$: the further below $k^*$, the
larger the excess of saving over depreciation per unit of capital.

\textbf{9.} $(1 - 0.35) \times 0.05 = 0.0325$; the half-life is about $\ln 2/0.0325 \approx 21$ years, close to the 23 found by iterating.

\textbf{10.} No: in the model, growth of output per worker slows to zero as $k$ approaches $k^*$. Lasting growth at 0.97\,\% a year needs rising technology.

\textbf{11.} $k^{*\prime} = (0.25 \times 10.166/0.05)^{1/0.65} \approx 421.5$; $y^{*\prime} \approx 84.3$.

\textbf{12.} $84.3/79.6 - 1 \approx 5.8\,\%$. Check: $(0.25/0.225)^{0.35/0.65} = 1.111^{0.538} \approx 1.058$.

\textbf{13.} Before: $0.775 \times 79.64 \approx 61.7$. After: $0.75 \times 84.29 \approx 63.2$. Consumption per worker is higher in the new \omterm{def:g12:solow-first-look:steady-state}{steady state}.

\textbf{14.} Zero: output per worker settles at the new, higher level.

\textbf{15.} No: with $s = 1$, consumption $(1 - s)y^*$ is zero, however large $y^*$. Raising the \omterm{def:g12:solow-first-look:saving-rate}{saving rate} raises consumption at first, then lowers it once the extra output
no longer covers the extra saving.

\textbf{16.} $k^* = (sA/\delta)^{1/(1 - \alpha)}$ rises with $A$, so the \omterm{def:g12:solow-first-look:steady-state}{steady state} moves up every year, and capital and output per worker keep growing towards a moving
target.

\textbf{17.} Fast convergence: Germany's output per person, below half its 1938 level in 1946, was above it by 1955 and half as high again by 1960. When skills,
institutions and technology survive, only capital needs rebuilding.

\textbf{18.} $k^*$ and $y^*$ rise with $A$: with a lower $A$ each unit of capital produces less, so less is saved, and saving matches depreciation at a lower $k$.

\textbf{19.} Growth of the workforce, which dilutes capital; technical progress; \omterm{def:g11:labour-market:human-capital}{human capital}; or a \omterm{def:g12:solow-first-look:saving-rate}{saving rate} that changes with income.

\textbf{20.} Brenland's steady-state output per worker is about 79.6 thousand 2015 dalens; after the \omterm{def:g12:solow-first-look:saving-rate}{saving rate} rises to 0.25, the model takes 22 years to close half the gap
to the new \omterm{def:g12:solow-first-look:steady-state}{steady state}.
\end{solution}
